\documentclass[12pt]{article}
\usepackage[T1]{fontenc}
\usepackage[utf8]{inputenc}
\usepackage{lmodern}
\usepackage{textcomp}
\usepackage{enumitem}
\usepackage{geometry}
\usepackage{graphicx}
\usepackage{amsmath, amssymb}
\geometry{margin=1in}
\begin{document}
\begin{center}
Economics - K-12 Level Assessment 22 \
Total Marks: 40 \
Answer all questions.
\end{center}

\section*{Multiple Choice (10 marks)}
\begin{enumerate}[label=\arabic*.]
\item Which of the following is NOT included in GDP?
\begin{enumerate}[label=(\alph*)]
    \item Federal government purchases of goods and services.
    \item Imports.
    \item State and local government purchases of goods and services.
    \item Exports.
\end{enumerate}
\item Which of the following indicates that two goods are complements?
\begin{enumerate}[label=(\alph*)]
    \item A positive income elasticity
    \item A horizontal demand curve
    \item A negative cross-price elasticity
    \item A demand elasticity greater than one
\end{enumerate}
\item An effective price ceiling in the market for good X likely results in
\begin{enumerate}[label=(\alph*)]
    \item a persistent surplus of good X.
    \item a persistent shortage of good X.
    \item an increase in the demand for good Y, a substitute for good X.
    \item a decrease in the demand for good Z, a complement with good X.
\end{enumerate}
\item Which of the following policies best describes supply-side fiscal policy?
\begin{enumerate}[label=(\alph*)]
    \item An increase in the money supply
    \item Increased government spending
    \item Lower taxes on research and development of new technology
    \item Higher taxes on household income
\end{enumerate}
\item The demand curve for a perfectly competitive firm's product is
\begin{enumerate}[label=(\alph*)]
    \item downward sloping and equal to the market demand curve.
    \item perfectly elastic.
    \item perfectly inelastic.
    \item kinked at the going market price.
\end{enumerate}
\end{enumerate}

\section*{True / False (10 marks)}
\textit{Answer True or False}
\begin{enumerate}[label=\arabic*.]
\setcounter{enumi}{5}
\item True or False: Marginal revenue equals marginal cost at the point where total revenue is maximized relative to total cost.
\item True or False: A change in the price of hamburgers does not shift the demand curve for hamburgers.
\item True or False: A perfectly competitive employer hires labor until the wage equals the marginal revenue product of labor.
\item True or False: An increase in the marginal propensity to save tends to increase the spending multiplier in an economy.
\item True or False: Households provide resources to firms in exchange for wages and profits.
\end{enumerate}

\section*{Long Form Response (20 marks)}
\textit{Answer each question with a clear and complete explanation.}
\begin{enumerate}[label=\arabic*.]
\setcounter{enumi}{10}
\item Discuss the implications of equilibrium in a perfectly competitive labor market and explain how it ensures that everyone who wants to work has the opportunity to do so.
\item Discuss the concept of diminishing marginal utility and explain how it relates to consumer behavior when the price of a good increases. Provide examples to support your analysis.
\end{enumerate}

\vfill
\noindent\textit{End of Paper}
\end{document}